% Lösungen Einheit 4 von 5 – Flächenbedingungen (zusammenbau v0.1)
\einheitenkopf{Einheit 4 von 5 · Flächenbedingungen}

\erg{15}{a) vorwärts (Inhalt gesucht) – Grenzen bestimmen, integrieren \quad b) Schnittstellen $0$ und $\frac{m}{2}$; Inhalt $\frac{m^3}{24} = 9$; $m = 6$ \quad c) Rechteck $6$ plus Dreieck $\frac{3}{2}(b - 2)$; $6 + \frac{3}{2}(b - 2) = 12$; $b = 6$ \quad d) Gesamt $8$; $k^3 = 4$; $k = \sqrt[3]{4} \approx 1{,}59$ \quad e) Nullstellen $\pm \frac{2}{\sqrt{a}}$; Inhalt $\frac{32}{3\sqrt{a}} = 16$; $a = \frac{4}{9} \approx 0{,}44$ \quad f) Der Flächenterm wächst mit $b$ stetig und streng monoton von null über alle Grenzen, weil $f$ positiv ist; also nimmt er den Wert fünf genau einmal an. \quad g) $k = -1$: Das Intervall von $-1$ bis $1$ liegt symmetrisch zu $0$; das Rechteck mit Höhe $3$ und Breite $2$ hat den Inhalt $6$, die Flächenstücke über und unter der Geraden $y = 3$ sind wegen der Punktsymmetrie gleich groß und heben sich auf.}
\erg{16}{a) Fehler: der Funktionswert statt der Fläche halbiert. Richtig: Gesamt $16$, $\frac{k^3}{4} = 8$, $k = \sqrt[3]{32} \approx 3{,}17$ \quad b) Stetig heißt ohne Sprung, er trifft also jeden Zwischenwert; streng monoton heißt, er erreicht keinen Wert ein zweites Mal. \quad c) Gesamt $\frac{40}{3}$; $\frac{k^3}{12} + 2k = \frac{20}{3}$ liefert $k \approx 2{,}60$; Abstand $\approx 26$ m}
